% ==============================================================================
% CAPÍTULO 2 - REVISÃO BIBLIOGRÁFICA
% ==============================================================================

\chapter{REVISÃO BIBLIOGRÁFICA} 
\label{ch:revisao-bibliografica}

Nesta seção, você deve fundamentar tecnicamente o seu estágio. 
Escreva um parágrafo introdutório apresentando os principais temas que 
serão discutidos e utilize subseções para organizar cada tecnologia ou 
conceito (ex: Linguagem Java, Banco de Dados, Frameworks). 
O objetivo é explicar as teorias e ferramentas utilizadas, demonstrando 
domínio técnico através de citações de autores ou documentações oficiais.


Abaixo, estão disponibilizados exemplos de como realizar as citações no texto utilizando os comandos do LaTeX.
Lembre-se de que todas as fontes citadas aqui devem estar cadastradas no seu arquivo de referências (.bib) 
para que sejam geradas automaticamente na lista final do trabalho.

\textbf{Exemplos de citações:}

Use o \texttt{\textbackslash citeonline\{\}} quando você quiser citar o nome do autor
 diretamente no seu texto (o LaTeX deixará apenas o ano entre parênteses).\\
 \textbf{Exemplo:} Segundo \citeonline{almeida2022}, a escolha correta de algoritmos de otimização é fundamental para melhorar o desempenho e a velocidade de resposta em sistemas distribuídos.\\

Use o \texttt{\textbackslash cite\{\}} quando fizer uma afirmação e quiser indicar a fonte apenas ao final (o LaTeX colocará o nome e o ano entre parênteses).\\
\textbf{Exemplo:} A utilização de sistemas distribuídos exige uma análise detalhada sobre como os algoritmos de otimização podem reduzir o tempo de processamento dos dados \cite{almeida2022}.

\textbf{Citações usadas para gerar as referências de exemplo:}\\
Estes exemplos são utilizados para mostrar como citar diferentes tipos de fontes são apresentados ao final do documento nas referências.\\
Verifique os campos necessários no arquivo "referencias.bib" para que elas apareçam corretamente nas Referências:  \\
Dissertação, \cite{almeida2022}\\
Trabalho publicado em Conferência, \cite{cbsoft2023}\\
Página web, \cite{ferreira2020}\\
Manual, \cite{latex2023}\\
Capitulo de livro, \cite{mendes2019}\\
Relatório técnico, \cite{oliveira2020}\\
Livro, \cite{linden2012}\\
Tese de doutorado, \cite{rodrigues2021}\\
Artigo científico, \cite{silva2023}\\
Livro de Conferência, \cite{souza2022}